\documentclass[11pt,a4paper]{article}

\usepackage{amsmath,amssymb,bm}
\usepackage{geometry}
\geometry{margin=1in}

\title{从协方差矩阵到 Cholesky 采样：一维高斯随机场离散采样笔记}
\author{}
\date{}

\begin{document}
\maketitle

\section{问题背景：离散点上的高斯随机场}

设一维空间区间为 $x\in[0,L]$，我们希望生成对数模量场
\[
\log E(x)
\]
作为高斯随机场（Gaussian Random Field, GRF）。在工程计算（例如 FEM 或 FNO 数据生成）中，我们只需要在离散网格点
\[
x_0,x_1,\dots,x_{n-1}
\]
上生成向量
\[
\bm{y}=
\begin{bmatrix}
\log E(x_0)\\
\log E(x_1)\\
\vdots\\
\log E(x_{n-1})
\end{bmatrix}.
\]

若假设 $\log E(x)$ 是高斯过程
\[
\log E(x)\sim \mathcal{GP}\!\left(\mu,\;k(x,x')\right),
\]
则离散后的向量 $\bm{y}$ 满足多元高斯分布
\[
\bm{y}\sim \mathcal{N}\!\left(\bm{\mu},\;\bm{K}\right),
\quad
\bm{\mu}=\mu\bm{1},
\quad
K_{ij}=k(x_i,x_j).
\]
其中 $\bm{K}$ 是协方差矩阵。

\section{协方差矩阵的含义（工程直觉）}

对任意两个位置 $x_i,x_j$，
\[
K_{ij}=\mathrm{Cov}\!\left[\log E(x_i),\log E(x_j)\right].
\]
因此：
\begin{itemize}
\item $K_{ii}=\mathrm{Var}[\log E(x_i)]$ 决定单点的波动强度（量级与 $\sigma^2$ 相关）。
\item $K_{ij}$（$i\neq j$）描述两点之间的相关性；距离越远，$K_{ij}$ 通常越小。
\end{itemize}

常见核函数（$\ell$ 为相关长度，$\sigma$ 为标准差尺度）：
\[
k_{\mathrm{exp}}(x,x')=\sigma^2\exp\!\left(-\frac{|x-x'|}{\ell}\right),
\qquad
k_{\mathrm{rbf}}(x,x')=\sigma^2\exp\!\left(-\frac{(x-x')^2}{2\ell^2}\right).
\]

\section{核心目标：如何生成 $\mathcal{N}(\bm{\mu},\bm{K})$ 的样本？}

已知一个结论：如果
\[
\bm{z}\sim \mathcal{N}(\bm{0},\bm{I}),
\]
并且存在矩阵 $\bm{A}$ 使得
\[
\bm{A}\bm{A}^\top=\bm{K},
\]
那么令
\[
\bm{y}=\bm{\mu}+\bm{A}\bm{z},
\]
就有
\[
\bm{y}\sim \mathcal{N}(\bm{\mu},\bm{K}).
\]

\subsection*{验证（为什么这个式子对？）}
由于 $\mathbb{E}[\bm{z}]=\bm{0}$，$\mathrm{Cov}(\bm{z})=\bm{I}$，
\[
\mathbb{E}[\bm{y}]=\mathbb{E}[\bm{\mu}+\bm{A}\bm{z}]=\bm{\mu},
\]
\[
\mathrm{Cov}(\bm{y})
=\mathrm{Cov}(\bm{A}\bm{z})
=\bm{A}\,\mathrm{Cov}(\bm{z})\,\bm{A}^\top
=\bm{A}\bm{I}\bm{A}^\top
=\bm{A}\bm{A}^\top
=\bm{K}.
\]
且线性变换保持高斯性，所以 $\bm{y}$ 仍为多元高斯。

\section{Cholesky 分解在这里做什么？}

对于对称正定（或数值上近似正定）的协方差矩阵 $\bm{K}$，Cholesky 分解给出
\[
\bm{K}=\bm{L}\bm{L}^\top,
\]
其中 $\bm{L}$ 是下三角矩阵。此时令 $\bm{A}=\bm{L}$ 即可。

所以采样流程变为：
\begin{enumerate}
\item 组装协方差矩阵：$K_{ij}=k(x_i,x_j)$；
\item 数值稳定：$\bm{K}\leftarrow \bm{K}+\varepsilon \bm{I}$（$\varepsilon$ 很小，例如 $10^{-8}$）；
\item Cholesky：$\bm{K}=\bm{L}\bm{L}^\top$；
\item 抽标准正态：$\bm{z}\sim \mathcal{N}(\bm{0},\bm{I})$；
\item 得到场样本：$\bm{y}=\bm{\mu}+\bm{L}\bm{z}$。
\end{enumerate}

\section{为何需要 ``jitter''：$\bm{K}+\varepsilon\bm{I}$？}

理论上由有效核函数构造的 $\bm{K}$ 应当是正定的，但在有限精度计算中可能出现：
\begin{itemize}
\item 网格点过密、$\ell$ 过大导致矩阵病态（条件数很大）；
\item 浮点误差使得最小特征值变成极小负数，从而 Cholesky 失败。
\end{itemize}
加入 $\varepsilon\bm{I}$ 等价于在每个点叠加一个极小的独立噪声方差，使矩阵数值上更稳定。

\section{与 ``白噪声'' 的关系（直觉图景）}

向量 $\bm{z}$ 是各分量独立的标准正态，工程上可理解为“白噪声”。左乘 $\bm{L}$ 相当于通过一个线性滤波器把白噪声变成具有空间相关性的“有色噪声”，从而得到相关长度受核函数控制的随机场样本。

\section{最终得到模量场}

采到 $\bm{y}=\log \bm{E}$ 后，取指数即可保证正性：
\[
E(x_i)=\exp(y_i).
\]
如果需要控制 $E$ 的均值或方差，通常在 $\log E$ 域通过 $\mu,\sigma$ 来设定更稳定。

\end{document}
